% This file was converted to LaTeX by Writer2LaTeX ver. 1.9.9
% see http://writer2latex.sourceforge.net for more info
\documentclass{article}
\usepackage{calc,amsmath,amssymb,amsfonts}
\usepackage[LGR,T1]{fontenc}
\usepackage[greek,italian]{babel}
\usepackage[style=numeric,backend=biber]{biblatex}
\author{HP}
\date{2026-07-26}
\begin{document}
\section{Prologo}
\subsection[Un pomeriggio d{}'estate]{Un pomeriggio d'estate}
Se chiudo gli occhi e provo a tornare indietro nel tempo, non rivedo subito un Commodore 64.

Rivedo il sole.

Quello delle estati interminabili, quando tre mesi di vacanza sembravano un'eternità e il tempo aveva un ritmo tutto
suo. Le giornate iniziavano presto, si interrompevano solo per il pranzo e riprendevano appena il caldo concedeva una
tregua. Bastava una bicicletta, un pallone o un'idea qualsiasi per trasformare un pomeriggio qualunque in un'avventura.

Era il 1982.

Avevo poco più di dieci anni e, come tanti bambini della mia età, non avevo la minima idea di quale sarebbe stato il mio
futuro.

C'era però una persona che, senza saperlo, stava per cambiarlo.

Si chiamava Emanuele.

Era il mio amico.

Anzi, per noi era l'amico tecnologico.

Suo padre gestiva un bar e preparava gelati artigianali. Emanuele lo aiutava spesso. Per noi bambini era quasi un lavoro
da grandi; per lui significava anche poter mettere da parte qualche soldo. All'epoca ci sembrava una piccola fortuna.

Fu così che, un giorno, a casa sua arrivò qualcosa che nessuno di noi aveva mai visto da vicino.

Una console per giocare a Pong.

Oggi sarebbe difficile spiegare l'effetto che ci fece. Abituati come siamo a immagini fotorealistiche, mondi
tridimensionali e giochi che sembrano film interattivi, due barre bianche e un quadratino che rimbalzava possono
sembrare poco più di un esercizio di fantasia.

Eppure, in quel momento, ci sembrò di avere il futuro davanti agli occhi.

Restavamo incantati.

Muovere una racchetta con una manopola e vedere una pallina reagire ai nostri movimenti era qualcosa di quasi magico.
Non era il gioco in sé ad affascinarci. Era il fatto che una macchina potesse rispondere alle nostre azioni.

Per la prima volta avevamo la sensazione di dialogare con un oggetto.

Non sapevamo ancora cosa fosse un microprocessore. Ignoravamo completamente il significato della parola software. E
probabilmente non avremmo saputo spiegare nemmeno cosa fosse davvero un computer.

Ma una domanda aveva già iniziato a farsi strada.

Come faceva quella macchina a sapere cosa doveva fare?

Fu una domanda destinata a non abbandonarmi più.

I pomeriggi da Emanuele non erano fatti soltanto di partite a Pong. Erano soprattutto fatti di racconti.

Qualcuno conosceva qualcuno che, pare, possedesse un VIC-20.

Si diceva che su quel computer girassero giochi incredibili. Con una grafica bellissima, suoni spettacolari e
possibilità che noi potevamo soltanto immaginare. Nessuno di noi li aveva mai visti davvero. Eppure ne parlavamo come
si parla di un luogo leggendario.

$\text{\textgreek{«}}$Pare che...$\text{\textgreek{»}}$

Era l'espressione con cui iniziavano tutte le storie.

Pare che il VIC-20 fosse un vero computer.

Pare che si potesse programmare.

Pare che non servisse soltanto per giocare.

Pare che si potessero inventare giochi nuovi.

Oggi basterebbero pochi secondi e una ricerca su Internet per verificare tutto. Allora, invece, le informazioni
viaggiavano da un amico all'altro, si arricchivano di particolari e alimentavano la fantasia. Ogni racconto sembrava
rendere quel misterioso computer ancora più straordinario.

Non avevo ancora capito che, senza accorgermene, stavo già cercando qualcosa di diverso.

Non volevo soltanto giocare.

Volevo capire.

Capire cosa si nascondesse dietro quei pixel.

Capire come una macchina potesse trasformare una serie di impulsi elettrici in un'avventura.

Capire, soprattutto, se un giorno anch'io sarei stato capace di costruire un videogioco.

Non immaginavo che la risposta sarebbe arrivata l'anno successivo, grazie a un professore che non avevo mai incontrato e
a un computer che portava un numero apparentemente insignificante.

Sessantaquattro.

Il nome continuava a ronzarmi in testa.

Sessantaquattro.

Non avevo la minima idea di cosa significasse quel numero. Non sapevo che si riferisse alla memoria del computer, non
conoscevo ancora il significato di un kilobyte e, soprattutto, non immaginavo che un giorno avrei saputo distinguere un
registro del processore da una locazione di memoria. Per me era semplicemente un nome. Eppure aveva qualcosa di diverso
da tutti gli altri.

Nei giorni successivi non si parlava d'altro. Ogni nuova informazione arrivava attraverso racconti di seconda mano,
quasi fossero leggende tramandate tra ragazzi. C'era sempre qualcuno che conosceva qualcuno che aveva visto quel
computer. Ogni volta il racconto si arricchiva di un particolare nuovo. Disegnava. Faceva i giochi. Si collegava al
televisore. Si potevano scrivere programmi. Qualcuno sosteneva addirittura che fosse capace di fare qualsiasi cosa,
purché gli si dicesse come.

Fu proprio quell'ultima affermazione a colpirmi più di tutte.

$\text{\textgreek{«}}$Purché gli si dicesse come.$\text{\textgreek{»}}$

Quelle quattro parole, senza che allora potessi comprenderlo, racchiudevano già l'essenza della programmazione.

Non era il computer a sapere cosa fare.

Era una persona a doverglielo insegnare.

Oggi questa osservazione può sembrare banale. Nel 1983, per un bambino che aveva visto solo videogiochi già pronti,
rappresentava una scoperta enorme. Per la prima volta intuivo che dietro ogni schermata, dietro ogni astronave che si
muoveva, dietro ogni pallina che rimbalzava, non c'era magia. C'era qualcuno che aveva scritto delle istruzioni.

Ed era proprio quel {\textquotedbl}qualcuno{\textquotedbl} che desideravo diventare.

Passarono poche settimane quando venni a sapere che nella scuola media del mio paese sarebbe stato organizzato un corso
pomeridiano di informatica. La notizia aveva però un piccolo problema: io frequentavo ancora la quinta elementare.

Il corso era destinato agli studenti delle medie.

Per qualche giorno rimasi combattuto. Da una parte pensavo che sarebbe stato inutile chiedere. Dall'altra continuavo a
ripetermi che quella poteva essere l'unica occasione per vedere davvero quel computer di cui tutti parlavano.

Alla fine decisi di provarci.

Ricordo perfettamente quel pranzo.

Eravamo seduti a tavola io, mia madre e mio padre. Non preparai nessun discorso. Non ce n'era bisogno. In casa avevano
già capito che il mio interesse per la tecnologia non era un capriccio destinato a durare una settimana. Ogni volta che
compariva qualcosa di nuovo, cercavo di capire come funzionasse. Smontavo, osservavo, facevo domande. Più che
utilizzare gli oggetti, volevo comprenderli.

Quando parlai del corso, mia madre ascoltò con attenzione. Mio padre, uomo di poche parole, non manifestò alcuna
contrarietà. Anzi, il fatto che non si opponesse era già una risposta.

La vera difficoltà era un'altra.

Convincere il preside ad accettare un bambino di quinta elementare in un corso pensato per ragazzi più grandi.

Fu mia madre a chiedere un colloquio.

Non ricordo grandi tensioni né particolari paure. Dentro di me speravo semplicemente che ci dessero una possibilità. In
fondo non si trattava di un'attività istituzionale della scuola, ma dell'iniziativa di un insegnante che aveva deciso
di condividere la propria passione con chiunque avesse voglia di imparare.

La risposta arrivò con una semplicità disarmante.

$\text{\textgreek{«}}$Va bene, proviamo.$\text{\textgreek{»}}$

All'epoca quelle tre parole mi sembrarono soltanto un'autorizzazione.

Oggi so che furono molto di più.

Furono il primo bivio della mia vita professionale.

Entrai in quell'aula con la curiosità di chi sta per vedere qualcosa di cui ha sentito parlare per settimane ma che non
ha mai avuto davanti agli occhi.

Era una normalissima aula di una scuola media dei primi anni Ottanta. La lavagna nera occupava la parete di fronte ai
banchi. La cattedra, consumata dal tempo, era posizionata come in qualsiasi altra classe. Se qualcuno fosse entrato in
quel momento, senza sapere cosa stava per accadere, avrebbe pensato di assistere a una comune lezione scolastica.

E invece, accanto alla cattedra, c'era lui.

Il Commodore 64.

Per la prima volta quel nome aveva finalmente assunto una forma.

La tastiera color beige, i tasti marroni, il logo con le strisce colorate nell'angolo superiore. Accanto, il monitor
Commodore. Poco più in là, un oggetto ancora più misterioso, con uno sportellino frontale e una leva metallica.

Era il drive.

Non avevo mai visto un lettore di floppy disk. Fino a quel momento i computer erano esistiti soltanto nei racconti degli
altri. Adesso erano lì, a pochi metri da me.

Non ricordo se in quell'istante osservassi più il computer o il professore.

Probabilmente entrambi.

Il professor Rossi Brunori non aveva l'atteggiamento di chi voleva dimostrare quanto fosse preparato. Al contrario,
sembrava divertirsi quanto noi. Parlava con calma, sorrideva spesso e aveva quella rara capacità di spiegare anche gli
argomenti più nuovi come se fossero assolutamente naturali.

Non c'era alcun bisogno di impressionare gli studenti.

Bisognava soltanto accendere la loro curiosità.

Oggi, dopo tanti anni trascorsi a progettare software, a seguire corsi di formazione e a spiegare informatica ad altre
persone, credo di aver capito quale fosse il suo vero talento.

Non insegnava il BASIC.

Insegnava a non averne paura.

È una differenza enorme.

Molti insegnanti iniziano mostrando quanto una materia sia complessa. Lui faceva l'esatto contrario. Ti convinceva che
fosse alla tua portata. Un problema dopo l'altro, una spiegazione dopo l'altra, costruiva lentamente la fiducia
necessaria per affrontare quello successivo.

Forse è proprio questo il motivo per cui ricordo così poco le singole lezioni e così tanto le sensazioni che provavo
uscendo dall'aula.

Ogni volta avevo imparato qualcosa che il giorno prima mi sembrava impossibile.

Il professore non utilizzava la lavagna.

Avrebbe potuto farlo, naturalmente. Era lì, alle sue spalle, pronta per essere riempita di schemi, frecce e formule. Ma
preferiva un altro approccio.

Scriveva direttamente sul Commodore.

Ogni comando compariva sullo schermo nel momento stesso in cui lo spiegava. Ogni errore diventava un'occasione per
capire meglio il funzionamento del linguaggio. Ogni modifica produceva immediatamente un risultato visibile.

Non c'era distanza tra la teoria e la pratica.

Coincidevano.

Noi eravamo tutti raccolti intorno a quella macchina. Non esisteva un computer per ogni studente. Ce n'era uno solo.
Eppure nessuno aveva la sensazione di partecipare passivamente.

Ogni programma che nasceva su quello schermo sembrava appartenere un po' a tutti.

Ricordo ancora il rumore dei tasti mentre il professore digitava il codice.

All'epoca era un suono come tanti.

Oggi, ripensandoci, mi sembra quasi il rumore con cui ha avuto inizio la mia professione.

Un particolare mi colpì fin dal primo giorno.

Il professore arrivava con il Commodore, il monitor e il drive trasportandoli personalmente. Ad aiutarlo c'era suo
figlio, più o meno della nostra età, che partecipava con naturalezza a quella piccola operazione di montaggio.
Sistemavano tutto con calma, collegavano i cavi e, pochi minuti dopo, quella normale aula scolastica si trasformava nel
luogo più moderno che avessi mai visto.

Non era un laboratorio di informatica.

Era molto di più.

Era un laboratorio di possibilità.

Quando tutto fu collegato, il professor Rossi Brunori si voltò verso di noi con la naturalezza di chi stava per iniziare
una normale lezione. Nessun gesto teatrale, nessuna pausa studiata per creare suspense. Eppure, senza che lui lo
sapesse, quel semplice movimento racchiudeva tutta l'attesa che avevo accumulato nelle settimane precedenti. Per la
prima volta avrei visto un Commodore 64 prendere vita davanti ai miei occhi.

Premette l'interruttore e, dopo un brevissimo istante, il monitor si illuminò. Lo schermo assunse quel caratteristico
colore azzurro chiaro incorniciato dal bordo blu, mentre comparvero alcune righe di testo che, allora, non
significavano assolutamente nulla per me. Erano parole e numeri che osservavo con la stessa curiosità con cui si guarda
un oggetto proveniente da un mondo sconosciuto. Non avevo ancora gli strumenti per comprenderli, ma intuivo che
ciascuna di quelle informazioni doveva avere un significato preciso.

Il professore lasciò che osservassimo lo schermo per qualche secondo. Non iniziò immediatamente a spiegare. Voleva che
ci abituassimo all'idea che quello non fosse un televisore acceso, ma una macchina pronta ad ascoltare ciò che avremmo
avuto da dirle.

Sul monitor compariva una scritta che occupava le prime righe dello schermo. Parlava di Commodore 64 BASIC V2, indicava
una quantità di memoria disponibile e terminava con una parola brevissima, seguita da un punto. READY.

Per noi era soltanto l'ultimo elemento visualizzato sullo schermo. Per il professore, invece, rappresentava l'inizio di
tutto.

Ci spiegò che il computer non stava mostrando un semplice messaggio di benvenuto. In realtà ci stava comunicando che
aveva terminato la propria fase di avvio, aveva controllato la memoria disponibile ed era pronto a ricevere istruzioni.
Quella parola, così semplice e apparentemente insignificante, racchiudeva un concetto destinato ad accompagnarmi per
tutta la mia vita professionale: un computer non prende iniziative, non immagina ciò che desideriamo fare e non
anticipa le nostre intenzioni. Attende istruzioni e le esegue con assoluta precisione, nel bene e nel male.

All'epoca ascoltai quella spiegazione con l'entusiasmo di un bambino che stava entrando in un territorio completamente
inesplorato. Oggi, dopo decenni trascorsi a scrivere software, mi rendo conto che il professor Rossi Brunori ci stava
trasmettendo molto più di una semplice nozione sul funzionamento del Commodore 64. Senza utilizzare termini complessi e
senza ricorrere a definizioni accademiche, ci stava introducendo a uno dei principi fondamentali dell'informatica: il
computer è uno strumento straordinariamente potente, ma la sua intelligenza coincide esattamente con la qualità delle
istruzioni che riceve.

Fu proprio partendo da quella semplice schermata iniziale che il professore iniziò a svelarci il significato delle
informazioni riportate sul monitor. Ci raccontò che il numero sessantaquattro, presente nel nome del computer, non era
stato scelto per ragioni commerciali o perché suonasse bene, ma indicava una caratteristica tecnica che, all'epoca,
rappresentava qualcosa di eccezionale. Per comprendere il motivo, però, era necessario fare un piccolo passo indietro e
capire che cosa fosse realmente la memoria di un computer e perché ogni programma, anche il più semplice, avesse
bisogno di uno spazio nel quale conservare dati e istruzioni.

Fu la prima volta che sentii pronunciare parole come byte e kilobyte. Non ne compresi subito il significato, ma ricordo
perfettamente la sensazione di trovarmi davanti a un mondo nel quale ogni risposta generava immediatamente una nuova
domanda. Più imparavo, più cresceva il desiderio di capire che cosa si nascondesse dietro quelle poche righe
visualizzate sullo schermo, perché intuivo che non erano state scritte per impressionare l'utente, ma per raccontargli
qualcosa sulla macchina che aveva davanti.

Senza rendermene conto, stavo già imparando una lezione che avrebbe accompagnato tutto il mio percorso professionale.
Ogni volta che un sistema mostra un'informazione, vale la pena chiedersi quale sia il motivo per cui è stata inserita e
quale storia tecnica stia cercando di raccontare. Quel giorno non stavo semplicemente osservando la schermata iniziale
di un Commodore 64; stavo imparando che, dietro ogni dettaglio apparentemente banale, esiste sempre una scelta
progettuale che merita di essere compresa.

Il professor Rossi Brunori lasciò trascorrere ancora qualche istante, osservando i nostri sguardi fissi sul monitor.
Sorrise, come se sapesse perfettamente quali domande ci stessero passando per la testa.

$\text{\textgreek{«}}$Allora...$\text{\textgreek{»}}$ disse appoggiandosi con naturalezza alla cattedra.
$\text{\textgreek{«}}$Chi di voi mi sa dire perché questo computer si chiama Commodore
Sessantaquattro?$\text{\textgreek{»}}$

Nell'aula calò il silenzio.

Qualcuno accennò un timido sorriso, altri abbassarono lo sguardo. Io stesso provai a immaginare una risposta, ma non
avevo la minima idea di quale potesse essere.

Il professore si guardò intorno soddisfatto.

$\text{\textgreek{«}}$Perfetto. È proprio la risposta che mi aspettavo.$\text{\textgreek{»}}$

Fece una breve pausa, non per creare suspense, ma quasi per darci il tempo di formulare un'altra ipotesi.

$\text{\textgreek{«}}$Sapete qual è la differenza tra usare una macchina e conoscerla davvero?$\text{\textgreek{»}}$

Anche questa volta nessuno rispose.

$\text{\textgreek{«}}$La maggior parte delle persone si accontenta di sapere che una macchina funziona. A noi, invece,
interessa capire perché funziona. È una differenza sottile, ma è quella che distingue chi utilizza un computer da chi,
un giorno, sarà capace di costruire un programma.$\text{\textgreek{»}}$

Quelle parole mi colpirono immediatamente. Ancora non sapevo scrivere una sola riga di BASIC, eppure il professore stava
già parlando di costruire programmi. Non sembrava affatto un obiettivo irraggiungibile. Lo diceva con tale naturalezza
che, per un attimo, mi sembrò quasi inevitabile che prima o poi ci sarei riuscito anch'io.

Il professore indicò con il dito il nome visualizzato sul monitor.

$\text{\textgreek{«}}$Quel sessantaquattro non è stato scelto perché suonava bene. Indica una caratteristica tecnica
molto importante di questo computer. Però, prima di arrivarci, dobbiamo chiarire una cosa. Secondo voi, che cos'è la
memoria?$\text{\textgreek{»}}$

Questa volta qualche mano si alzò.

$\text{\textgreek{«}}$È dove il computer conserva i programmi.$\text{\textgreek{»}}$

$\text{\textgreek{«}}$Bravo. E anche i dati,$\text{\textgreek{»}}$ aggiunse il professore.
$\text{\textgreek{«}}$Immaginate la memoria come una lunga fila di piccoli cassetti. Ogni cassetto può contenere una
sola informazione e ognuno ha il proprio indirizzo. Quando il processore deve leggere o scrivere qualcosa, non cerca a
caso. Sa esattamente in quale cassetto andare.$\text{\textgreek{»}}$

Prese un gessetto dalla vaschetta della lavagna, lo fece ruotare tra le dita e poi sorrise.

$\text{\textgreek{«}}$No, tranquilli... oggi non riempiremo la lavagna di numeri.$\text{\textgreek{»}}$

Qualcuno rise.

$\text{\textgreek{«}}$Preferisco che il Commodore parli direttamente con voi.$\text{\textgreek{»}}$

Posò il gessetto senza aver scritto nulla e tornò davanti alla tastiera.

$\text{\textgreek{«}}$Ogni cassetto della memoria può contenere un byte. È una parola inglese che incontrerete spesso e
che, con il passare del tempo, vi diventerà familiare quanto il vostro nome. Oggi vi basta sapere che un byte
rappresenta la più piccola quantità di memoria con cui questo computer può lavorare.$\text{\textgreek{»}}$

Indicò nuovamente la scritta sul monitor.

$\text{\textgreek{«}}$Adesso facciamo un piccolo conto insieme. Se questo computer si chiama Commodore Sessantaquattro,
secondo voi quanti kilobyte di memoria avrà?$\text{\textgreek{»}}$

Questa volta le risposte arrivarono quasi tutte insieme.

$\text{\textgreek{«}}$Sessantaquattro!$\text{\textgreek{»}}$

$\text{\textgreek{«}}$Esatto.$\text{\textgreek{»}}$ Il professore annuì soddisfatto. $\text{\textgreek{«}}$E allora
osservate bene la schermata. Se possiede sessantaquattro kilobyte, perché il computer vi dice che i byte disponibili
sono trentottomilanovecentoundici? Gli altri dove sono finiti? Qualcuno li ha rubati durante la
notte?$\text{\textgreek{»}}$

L'aula scoppiò a ridere.

Io no.

Continuavo a fissare quel numero.

Per la prima volta mi accorsi che un computer non mostrava mai informazioni a caso. Se compariva quel valore,
significava che doveva esistere una spiegazione. E, ancora una volta, il professore non ci aveva dato una risposta. Ci
aveva regalato una domanda.

Col tempo avrei capito che era quello il suo modo di insegnare. Non cercava studenti capaci di memorizzare delle
nozioni; cercava ragazzi che imparassero a non smettere mai di farsi domande.

Il professor Rossi Brunori lasciò che la nostra curiosità crescesse ancora per qualche istante. Non aveva alcuna fretta
di rispondere. Era evidente che quel silenzio faceva parte della lezione tanto quanto le parole che avrebbe pronunciato
di lì a poco.

$\text{\textgreek{«}}$Vedete...$\text{\textgreek{»}}$ riprese con il tono tranquillo che ormai avevamo imparato a
conoscere, $\text{\textgreek{«}}$questa è una delle prime cose che bisogna imparare quando ci si avvicina a un
computer. I numeri non sono lì per caso. Se il computer vi mostra un'informazione, significa che sta cercando di dirvi
qualcosa. Sta a noi capire che cosa.$\text{\textgreek{»}}$

Si avvicinò nuovamente al monitor e indicò con un dito la riga che riportava la quantità di memoria disponibile.

$\text{\textgreek{«}}$Prima abbiamo detto che il Commodore 64 possiede sessantaquattro kilobyte di memoria. È vero. Però
non significa che possiamo utilizzarli tutti per scrivere i nostri programmi. Sarebbe un po' come acquistare una casa e
pensare di poter occupare ogni centimetro disponibile con i mobili, dimenticando che devono pur esserci le porte, le
finestre, il bagno e la cucina.$\text{\textgreek{»}}$

Qualcuno sorrise.

$\text{\textgreek{«}}$La memoria del computer funziona in modo molto simile. Una parte serve a voi, ma un'altra parte
serve al computer stesso per poter lavorare. Se la occupassimo tutta con i nostri programmi, il Commodore non saprebbe
più nemmeno come accendere lo schermo o interpretare quello che scrivete sulla tastiera.$\text{\textgreek{»}}$

Quelle parole resero immediatamente più comprensibile qualcosa che fino a pochi secondi prima sembrava astratto.

Il professore si sedette sul bordo della cattedra, quasi volesse trasformare quella spiegazione in una conversazione.

$\text{\textgreek{«}}$Immaginate di dover organizzare una grande biblioteca. Potreste riempire ogni scaffale di libri,
certo. Però vi servirebbe anche un catalogo, qualcuno che sappia dove sono sistemati i volumi e qualche corridoio per
poterci camminare in mezzo. Se eliminaste tutto questo per fare posto ad altri libri, la biblioteca diventerebbe
inutilizzabile.$\text{\textgreek{»}}$

Fece una breve pausa per osservare le nostre espressioni.

$\text{\textgreek{«}}$Il Commodore fa esattamente la stessa cosa. Tiene una parte della memoria per sé, perché gli serve
per funzionare. Quello che rimane lo mette a disposizione dei vostri programmi. Ed è proprio quella quantità che vedete
scritta sullo schermo.$\text{\textgreek{»}}$

Un ragazzo seduto nelle prime file alzò la mano.

$\text{\textgreek{«}}$Professore... allora i kilobyte mancanti non sono sprecati?$\text{\textgreek{»}}$

Rossi Brunori sorrise con soddisfazione.

$\text{\textgreek{«}}$Bella domanda. No, non sono affatto sprecati. Anzi, stanno lavorando anche quando voi non ve ne
accorgete. Pensate a quello che succede mentre io scrivo sulla tastiera. Ogni lettera compare immediatamente sullo
schermo. Se premo il tasto Cancella, il cursore torna indietro. Se sbaglio una parola, il computer sa come correggerla.
Tutto questo non accade per magia. Da qualche parte esiste già un programma che sa come farlo.$\text{\textgreek{»}}$

Si voltò verso il Commodore e appoggiò delicatamente una mano sulla tastiera.

$\text{\textgreek{«}}$Quel programma è già dentro il computer. Non dobbiamo caricarlo da una cassetta né da un floppy
disk. Fa parte della macchina stessa. È uno dei motivi per cui, appena lo accendiamo, possiamo iniziare immediatamente
a scrivere.$\text{\textgreek{»}}$

Quella spiegazione mi colpì profondamente.

Fino a quel momento avevo sempre immaginato il computer come un contenitore vuoto che si limitava a eseguire ciò che
qualcuno gli inseriva dall'esterno. Per la prima volta capii che possedeva già una sorta di patrimonio iniziale, un
insieme di istruzioni che gli permettevano di prendere vita ancora prima che l'utente digitasse il primo comando.

Il professore sembrò intuire che quel concetto aveva acceso qualcosa dentro di noi e decise di spingersi un po' oltre.

$\text{\textgreek{«}}$Adesso osservate quest'ultima parola.$\text{\textgreek{»}}$

Indicò il cursore lampeggiante accanto alla scritta READY.

$\text{\textgreek{«}}$Che cosa significa secondo voi?$\text{\textgreek{»}}$

Le risposte furono le più disparate.

$\text{\textgreek{«}}$Che è acceso.$\text{\textgreek{»}}$

$\text{\textgreek{«}}$Che funziona.$\text{\textgreek{»}}$

$\text{\textgreek{«}}$Che possiamo usarlo.$\text{\textgreek{»}}$

$\text{\textgreek{«}}$Sì...$\text{\textgreek{»}}$ rispose lui annuendo. $\text{\textgreek{«}}$Ma c'è qualcosa di ancora
più importante.$\text{\textgreek{»}}$

Lasciò trascorrere qualche secondo, poi continuò.

$\text{\textgreek{«}}$Ready significa {\textquotedbl}sono pronto{\textquotedbl}. È il computer che sta parlando con voi.
Vi sta dicendo che ha terminato tutto quello che doveva fare da solo e che, da questo momento in poi, aspetta le vostre
istruzioni. Finché non sarete voi a prendere un'iniziativa, lui rimarrà lì, pazientemente, ad
attendere.$\text{\textgreek{»}}$

Sul monitor il cursore continuava a lampeggiare con regolarità, quasi volesse confermare quelle parole.

$\text{\textgreek{«}}$Ed è proprio questo il bello della programmazione,$\text{\textgreek{»}}$ concluse il professore.
$\text{\textgreek{«}}$Ogni volta che vedrete comparire quella parola, ricordatevi che il computer vi sta consegnando il
controllo. Da quel momento in avanti, ciò che accadrà dipenderà dalle istruzioni che deciderete di
scrivere.$\text{\textgreek{»}}$

In quell'istante nessuno di noi poteva immaginare che quelle istruzioni sarebbero state, per molti, soltanto un
esercizio scolastico destinato a essere dimenticato nel giro di qualche mese. Per me, invece, rappresentavano l'inizio
di un viaggio che, molti anni dopo, mi avrebbe portato a progettare software, a vivere di informatica e, soprattutto, a
tornare su quei banchi con il desiderio di completare il videogioco che allora esisteva soltanto nella mia
immaginazione.

Il professor Rossi Brunori rimase qualche istante in silenzio, osservando il cursore che continuava a lampeggiare sullo
schermo. Non sembrava interessato a proseguire immediatamente la spiegazione. Al contrario, dava l'impressione di voler
lasciare che quella semplice parola, READY., si depositasse lentamente nella nostra mente.

Poi si voltò verso di noi.

$\text{\textgreek{«}}$Adesso vi faccio una domanda che, forse, vi sembrerà banale.$\text{\textgreek{»}}$

Sorrise, come faceva ogni volta che stava per introdurre un concetto importante.

$\text{\textgreek{«}}$Secondo voi, come si parla con un computer?$\text{\textgreek{»}}$

La domanda ci colse impreparati.

Eravamo abituati a pensare al computer come a una macchina straordinaria, ma nessuno si era mai posto il problema di
come comunicare con lui. Per qualche secondo si sentì soltanto il leggero ronzio del monitor.

$\text{\textgreek{«}}$Con la tastiera...$\text{\textgreek{»}}$ rispose finalmente uno dei ragazzi.

$\text{\textgreek{«}}$È vero,$\text{\textgreek{»}}$ confermò il professore. $\text{\textgreek{«}}$La tastiera è lo
strumento che utilizziamo. Ma io non vi ho chiesto con che cosa si parla al computer. Vi ho chiesto
come.$\text{\textgreek{»}}$

La differenza tra le due domande ci sfuggì completamente.

Il professore non mostrò alcun segno di impazienza. Anzi, sembrava soddisfatto che nessuno avesse ancora trovato la
risposta.

$\text{\textgreek{«}}$Provate a immaginare di andare in un Paese del quale non conoscete la lingua. Avete una bocca,
avete una voce, potete parlare quanto volete, ma se utilizzate parole che nessuno comprende, riuscirete davvero a
comunicare?$\text{\textgreek{»}}$

Questa volta le teste iniziarono ad annuire.

$\text{\textgreek{«}}$No.$\text{\textgreek{»}}$

$\text{\textgreek{«}}$Esatto. Avere una tastiera non basta, proprio come avere una bocca non basta per parlare con una
persona. Per comunicare è necessario condividere un linguaggio.$\text{\textgreek{»}}$

Si avvicinò lentamente al Commodore e appoggiò le mani sulla tastiera.

$\text{\textgreek{«}}$Questo computer non conosce l'italiano. Non conosce nemmeno l'inglese, almeno non nel modo in cui
lo intendiamo noi. Conosce un linguaggio molto più semplice, costruito apposta perché possa comprenderlo senza
ambiguità.$\text{\textgreek{»}}$

Mentre parlava, accarezzava con le dita alcuni tasti, quasi fossero le pagine di un libro che aveva imparato a conoscere
molti anni prima.

$\text{\textgreek{«}}$Quel linguaggio si chiama BASIC.$\text{\textgreek{»}}$

Pronunciò quella parola con una naturalezza disarmante, senza attribuirle alcun alone di mistero.

$\text{\textgreek{«}}$Non lasciatevi spaventare dal nome. Tra qualche settimana vi sembrerà normale quanto scrivere una
frase su un quaderno. La differenza è che, invece di rivolgervi a un insegnante o a un compagno di classe, starete
parlando con una macchina.$\text{\textgreek{»}}$

Si interruppe per osservare le nostre reazioni.

$\text{\textgreek{«}}$Naturalmente il computer è molto esigente. Se dimenticate una lettera, se invertite una parola o
se scrivete qualcosa che non appartiene al suo linguaggio, lui non proverà a indovinare quello che volevate dire. Si
limiterà a fermarsi e vi dirà che non ha capito.$\text{\textgreek{»}}$

Rise leggermente.

$\text{\textgreek{«}}$Qualcuno potrebbe pensare che sia poco intelligente. Io preferisco dire che è incredibilmente
sincero.$\text{\textgreek{»}}$

Nell'aula si levò qualche risata.

$\text{\textgreek{«}}$Le persone, qualche volta, fingono di aver capito per non mettervi in imbarazzo. Un computer non
lo farà mai. Se non comprende un'istruzione, ve lo dirà immediatamente. Può sembrare un difetto, ma in realtà è uno dei
suoi più grandi pregi, perché vi costringerà a ragionare con precisione.$\text{\textgreek{»}}$

Quelle parole, allora, sembravano soltanto un'osservazione curiosa. Ripensandoci oggi, mi accorgo che contenevano una
lezione molto più profonda, una lezione che continua a valere anche per i linguaggi di programmazione moderni. Il
computer non è severo per cattiveria; è coerente. E proprio questa coerenza rende possibile costruire programmi sempre
più complessi, sapendo che ogni istruzione verrà interpretata esattamente nello stesso modo.

Il professore si voltò nuovamente verso il monitor.

$\text{\textgreek{«}}$Adesso basta parlare di linguaggi. Vediamone uno all'opera.$\text{\textgreek{»}}$

Le sue mani iniziarono a muoversi sulla tastiera con una sicurezza che mi lasciò affascinato. Non cercava i tasti, non
abbassava lo sguardo, non sembrava dover riflettere su ciò che stava facendo. Era evidente che quel computer fosse
ormai un compagno di lavoro abituale, quasi un'estensione naturale del suo modo di pensare.

Sul monitor comparvero lentamente due caratteri.

10

Il professore si fermò.

$\text{\textgreek{«}}$No, non ho sbagliato a premere Invio,$\text{\textgreek{»}}$ disse anticipando la domanda che
probabilmente tutti ci stavamo facendo. $\text{\textgreek{«}}$La prima cosa che si scrive in un programma BASIC non è
un comando, ma un numero.$\text{\textgreek{»}}$

Si voltò verso di noi con un sorriso divertito.

$\text{\textgreek{«}}$E adesso immagino che vi starete chiedendo perché mai un programma dovrebbe cominciare con un
numero invece che con una parola.$\text{\textgreek{»}}$

Non aspettò alcuna risposta.

$\text{\textgreek{«}}$È un'ottima domanda. E, come succede spesso in informatica, la risposta è molto più interessante
di quanto possiate immaginare.$\text{\textgreek{»}}$

Fece una breve pausa, poi aggiunse con tono quasi confidenziale:

$\text{\textgreek{«}}$Se oggi comprenderete il motivo per cui un programma BASIC inizia con un numero, avrete già
imparato qualcosa che molti utenti di un computer non si chiedono nemmeno dopo anni di utilizzo. E noi non siamo qui
per imparare a usare un computer. Siamo qui per capire come ragiona.$\text{\textgreek{»}}$

Il professor Rossi Brunori rimase ancora qualche secondo davanti alla tastiera. Era evidente che sapesse perfettamente
quale sarebbe stato il passo successivo, ma sembrava volerci accompagnare fino a quel punto senza mai farci sentire in
ritardo.

$\text{\textgreek{«}}$Bene,$\text{\textgreek{»}}$ disse infine. $\text{\textgreek{«}}$Adesso il Commodore è pronto ad
ascoltarci. Tocca a noi dirgli qualcosa.$\text{\textgreek{»}}$

Si voltò verso la classe.

$\text{\textgreek{«}}$Secondo voi, qual è la prima cosa che si dice a una persona quando la si
incontra?$\text{\textgreek{»}}$

$\text{\textgreek{«}}$Ciao!$\text{\textgreek{»}}$ rispose qualcuno dalle ultime file.

$\text{\textgreek{«}}$Esatto.$\text{\textgreek{»}}$ Il professore sorrise. $\text{\textgreek{«}}$È un'ottima idea.
Allora proviamo a salutare anche il computer.$\text{\textgreek{»}}$

Tornò alla tastiera e, accanto al numero che aveva già scritto, iniziò a digitare lentamente.

10 PRINT {\textquotedbl}CIAO{\textquotedbl}

Non premette subito il tasto Invio.

Lasciò che osservassimo quella riga.

$\text{\textgreek{«}}$Leggetela con me.$\text{\textgreek{»}}$

Lo facemmo quasi in coro.

$\text{\textgreek{«}}$Dieci... Print... Ciao.$\text{\textgreek{»}}$

Il professore scosse leggermente la testa.

$\text{\textgreek{«}}$No. Se volete imparare davvero a programmare, dovete abituarvi a leggere quello che il computer
legge, non quello che vede un italiano.$\text{\textgreek{»}}$

Ripeté indicando con il dito ogni parola.

$\text{\textgreek{«}}$Linea dieci. Comando PRINT. Testo CIAO.$\text{\textgreek{»}}$

Ripeté la frase una seconda volta.

$\text{\textgreek{«}}$Linea dieci. Comando PRINT. Testo CIAO.$\text{\textgreek{»}}$

Quella semplice correzione poteva sembrare insignificante, ma cambiava completamente il modo di osservare il programma.
Non stavamo più leggendo una frase; stavamo analizzando tre elementi diversi, ognuno con un significato preciso.

$\text{\textgreek{«}}$Sapete perché abbiamo scritto il numero dieci?$\text{\textgreek{»}}$

Qualcuno azzardò una risposta.

$\text{\textgreek{«}}$Perché è la prima riga.$\text{\textgreek{»}}$

$\text{\textgreek{«}}$Sì, ma avremmo potuto scrivere anche uno, oppure cento, oppure mille. Il computer avrebbe eseguito
comunque il programma.$\text{\textgreek{»}}$

Prese una pausa.

$\text{\textgreek{«}}$Allora perché proprio dieci?$\text{\textgreek{»}}$

Nessuno riuscì a trovare una spiegazione.

$\text{\textgreek{«}}$Perché chi programma deve pensare anche al domani.$\text{\textgreek{»}}$

Quella risposta ci lasciò ancora più confusi.

Il professore rise di gusto.

$\text{\textgreek{«}}$Lo so. In questo momento vi sembra una risposta senza senso. Però immaginate di tornare domani su
questo programma e di voler aggiungere un'istruzione tra la prima e la seconda riga. Se oggi numerassimo le righe uno,
due, tre e quattro, saremmo costretti a riscrivere tutto. Lasciando invece uno spazio di dieci in dieci, potremo
inserire comodamente una riga quindici, oppure una riga dodici, senza modificare il resto del
programma.$\text{\textgreek{»}}$

Si fermò un istante.

$\text{\textgreek{«}}$Programmare non significa soltanto scrivere istruzioni. Significa prevedere che, prima o poi,
qualcuno dovrà modificarle. Quel qualcuno, molto spesso, sarete proprio voi.$\text{\textgreek{»}}$

Ancora oggi, ripensando a quel momento, mi sorprende quanto fosse moderna quella osservazione. Il professore non ci
stava semplicemente spiegando una convenzione del BASIC. Senza utilizzare parole come \emph{manutenibilità},
\emph{evoluzione del software} o \emph{refactoring}, ci stava insegnando che un programma non nasce mai definitivo. Un
buon programmatore non pensa soltanto a far funzionare il codice oggi; cerca di renderlo comprensibile e modificabile
anche domani.

Il professor Rossi Brunori premette infine il tasto \textbf{RETURN}.

La riga scomparve e il cursore si posizionò immediatamente sulla linea successiva.

Il computer aveva accettato l'istruzione.

Non aveva ancora eseguito nulla.

L'aveva semplicemente memorizzata.

Fu allora che il professore si voltò verso di noi con un'espressione quasi divertita.

$\text{\textgreek{«}}$Ecco un'altra cosa che vale la pena ricordare. Molti pensano che un computer esegua tutto quello
che si scrive nel momento stesso in cui lo si digita. Il BASIC, invece, quando una riga inizia con un numero, si
comporta diversamente. Non la esegue. La conserva in memoria, aspettando che siate voi a dirgli quando
iniziare.$\text{\textgreek{»}}$

Quelle parole aprivano un'altra finestra su un mondo che non avevamo mai immaginato.

Il Commodore non era soltanto capace di eseguire istruzioni. Era anche capace di ricordarle, ordinarle e conservarle
fino al momento opportuno.

Era come se, lentamente, quella macchina stesse smettendo di essere un oggetto misterioso per trasformarsi in un
interlocutore con un proprio modo di ragionare.

Io, seduto tra quei banchi, non riuscivo più a staccare gli occhi dallo schermo.

Per qualche istante nell'aula non accadde assolutamente nulla.

Sul monitor rimase visibile il cursore lampeggiante, ora posizionato sulla riga successiva, come se il Commodore fosse
semplicemente in attesa di ricevere un'altra istruzione. Il saluto che avevamo appena scritto non compariva da nessuna
parte. Non c'era alcun {\textquotedbl}CIAO{\textquotedbl} al centro dello schermo, nessun effetto speciale, nessun
segnale che lasciasse intuire di aver ottenuto un risultato.

Il professor Rossi Brunori continuò a osservare il monitor senza dire una parola. Era evidente che stesse aspettando
proprio quella reazione.

Le espressioni dei ragazzi cambiarono lentamente. Alcuni si guardarono tra loro, altri tornarono a fissare la tastiera,
come se cercassero di capire quale passaggio fosse sfuggito. Anch'io provavo una leggera sensazione di delusione. Ero
convinto che, premendo il tasto \textbf{RETURN}, il computer avrebbe immediatamente eseguito il comando che avevamo
scritto.

Fu un ragazzo seduto nelle prime file a rompere il silenzio.

$\text{\textgreek{«}}$Professore... ma non funziona?$\text{\textgreek{»}}$

Rossi Brunori sorrise con l'aria di chi aveva appena sentito la domanda che sperava di ascoltare.

$\text{\textgreek{«}}$Secondo voi, perché dovrebbe funzionare?$\text{\textgreek{»}}$

La domanda ci spiazzò.

$\text{\textgreek{«}}$Perché glielo abbiamo scritto...$\text{\textgreek{»}}$

$\text{\textgreek{«}}$Davvero?$\text{\textgreek{»}}$

Il professore incrociò le braccia e si appoggiò leggermente alla cattedra.

$\text{\textgreek{«}}$Proviamo a ragionare insieme. Se io vi consegnassi un foglio con scritto il programma della
giornata di domani, significherebbe forse che dovete iniziare immediatamente a fare tutto quello che c'è
scritto?$\text{\textgreek{»}}$

$\text{\textgreek{«}}$No.$\text{\textgreek{»}}$

$\text{\textgreek{«}}$Perché no?$\text{\textgreek{»}}$

$\text{\textgreek{«}}$Perché è solo un programma.$\text{\textgreek{»}}$

$\text{\textgreek{«}}$Esattamente.$\text{\textgreek{»}}$

Il professore annuì soddisfatto.

$\text{\textgreek{«}}$Quella riga che abbiamo scritto non è un ordine eseguito. È un programma. Il Commodore l'ha letta,
l'ha capita e l'ha conservata nella sua memoria. Adesso sa che cosa dovrà fare, ma aspetta ancora che qualcuno gli dica
quando cominciare.$\text{\textgreek{»}}$

Si avvicinò nuovamente alla tastiera.

$\text{\textgreek{«}}$Vedete, questo è uno degli aspetti più intelligenti del BASIC. Voi potete scrivere un programma
lungo una riga, dieci righe o mille righe senza che il computer esegua nulla. Potete correggerlo, aggiungere nuove
istruzioni, eliminarne alcune e modificarlo tutte le volte che volete. Soltanto quando sarete soddisfatti gli direte:
{\textquotedbl}Adesso eseguilo{\textquotedbl}.$\text{\textgreek{»}}$

Mentre parlava, mi accorgevo che il Commodore stava assumendo una personalità completamente diversa da quella che avevo
immaginato. Non era una macchina impaziente che reagiva istantaneamente a ogni tasto premuto. Era sorprendentemente
disciplinata. Ascoltava, prendeva nota, aspettava. Soltanto dopo aver ricevuto un comando preciso iniziava a lavorare.

Il professor Rossi Brunori tornò a rivolgersi alla classe.

$\text{\textgreek{«}}$Secondo voi, come potremmo dire al computer che è arrivato il momento di eseguire il
programma?$\text{\textgreek{»}}$

Questa volta nessuno tentò di indovinare. Avevamo imparato che, dietro ogni domanda del professore, si nascondeva un
ragionamento che valeva la pena seguire fino in fondo.

Lui sorrise.

$\text{\textgreek{«}}$Il BASIC utilizza una parola molto semplice. Una parola inglese di appena tre
lettere.$\text{\textgreek{»}}$

Lasciò trascorrere qualche secondo, quasi volesse dare a ciascuno di noi la possibilità di formularla mentalmente.

$\text{\textgreek{«}}$RUN.$\text{\textgreek{»}}$

La digitò lentamente sulla tastiera, senza alcun numero davanti.

Poi si voltò verso di noi.

$\text{\textgreek{«}}$Vi ricordate quello che vi ho detto poco fa? Quando una riga inizia con un numero, il Commodore la
memorizza come parte del programma. Quando invece scrivete un comando senza numero, il computer lo esegue
immediatamente. È una distinzione fondamentale, perché da questo momento in poi dovrete sempre sapere se state
costruendo un programma oppure se state dialogando direttamente con la macchina.$\text{\textgreek{»}}$

Per un istante nessuno parlò.

Avevo la netta sensazione che stesse per accadere qualcosa di importante. Quel piccolo programma, composto da una sola
istruzione, era ancora fermo nella memoria del Commodore, in attesa di ricevere il permesso di vivere.

Il professore appoggiò delicatamente l'indice sul tasto \textbf{RETURN}, ma non lo premette subito.

$\text{\textgreek{«}}$Prima di eseguirlo,$\text{\textgreek{»}}$ disse con un sorriso, $\text{\textgreek{«}}$proviamo a
fare una previsione. Un bravo programmatore non aspetta che sia il computer a dirgli cosa succederà. Cerca di
immaginarlo prima. Allora ditemi: quando premerò questo tasto, che cosa vi aspettate di vedere sullo
schermo?$\text{\textgreek{»}}$

Quella domanda rimase sospesa nell'aula, insieme all'attesa di un gruppo di ragazzi che, forse per la prima volta, non
stavano semplicemente osservando un computer, ma iniziavano davvero a ragionare come programmatori.

Per qualche secondo nessuno ebbe il coraggio di rispondere.

Non era timidezza. Era qualcosa di diverso. Avevamo imparato che le domande del professor Rossi Brunori non erano mai
poste per verificare se conoscessimo una risposta, ma per insegnarci a costruirla. Per questo motivo ognuno di noi
cercava di immaginare ciò che sarebbe accaduto, mettendo insieme tutto quello che aveva ascoltato fino a quel momento.

Fu il ragazzo seduto vicino alla finestra a rompere il silenzio.

$\text{\textgreek{«}}$Secondo me scriverà {\textquotedbl}CIAO{\textquotedbl}.$\text{\textgreek{»}}$

Il professore annuì.

$\text{\textgreek{«}}$Bene. Dove lo scriverà?$\text{\textgreek{»}}$

Il ragazzo esitò.

$\text{\textgreek{«}}$Al centro dello schermo... forse.$\text{\textgreek{»}}$

Un altro intervenne quasi subito.

$\text{\textgreek{«}}$Io penso sotto il programma.$\text{\textgreek{»}}$

$\text{\textgreek{«}}$Perché?$\text{\textgreek{»}}$ domandò il professore.

$\text{\textgreek{«}}$Perché la riga è rimasta in memoria.$\text{\textgreek{»}}$

$\text{\textgreek{«}}$Interessante.$\text{\textgreek{»}}$

Non disse altro. Non fece capire quale delle due ipotesi fosse corretta. Continuava a raccogliere le nostre idee come se
ciascuna meritasse di essere ascoltata.

Poi il suo sguardo si fermò su di me.

$\text{\textgreek{«}}$Enzo, tu che cosa ne pensi?$\text{\textgreek{»}}$

Sentii gli occhi di tutta la classe voltarsi nella mia direzione.

Abbassai lo sguardo verso il banco, poi tornai a osservare il monitor.

$\text{\textgreek{«}}$Non lo so...$\text{\textgreek{»}}$ risposi sinceramente. $\text{\textgreek{«}}$Però penso che farà
esattamente quello che gli abbiamo scritto.$\text{\textgreek{»}}$

Il professore sorrise.

$\text{\textgreek{«}}$Questa è una risposta che mi piace.$\text{\textgreek{»}}$

Rimase qualche istante in silenzio, poi aggiunse con voce pacata:

$\text{\textgreek{«}}$Non perché sia completa, ma perché contiene l'idea più importante di tutte. Un computer non fa
quello che desideriamo. Fa quello che gli chiediamo di fare. E il nostro compito, da oggi in avanti, sarà imparare a
chiederglielo nel modo giusto.$\text{\textgreek{»}}$

Quelle parole, all'epoca, mi sembrarono semplicemente una buona spiegazione. Molti anni dopo mi sarei accorto che
rappresentavano una delle definizioni più semplici e, allo stesso tempo, più accurate della programmazione che abbia
mai ascoltato.

Il professor Rossi Brunori tornò verso la tastiera.

Sul monitor comparvero tre sole lettere.

RUN

Per un attimo nessuno respirò.

Il tasto \textbf{RETURN} venne premuto con la stessa naturalezza con cui, pochi minuti prima, era stato acceso il
computer.

Lo schermo sembrò fermarsi per una frazione di secondo.

Poi accadde.

Sotto il comando apparve una sola parola.

CIAO

Nient'altro.

Nessuna musica.

Nessuna animazione.

Nessun effetto speciale.

Soltanto cinque lettere bianche su uno sfondo azzurro.

Eppure, in quell'istante, mi sembrarono la cosa più straordinaria che avessi mai visto.

Non stavo osservando una dimostrazione preparata in anticipo. Quelle lettere non erano già presenti nella memoria del
Commodore. Erano comparse perché qualcuno aveva scritto un'istruzione e il computer l'aveva eseguita esattamente come
richiesto.

Fu in quel preciso momento che iniziai a comprendere quale fosse la vera natura di quella macchina.

Non era il computer a essere interessante.

Interessanti erano le istruzioni.

Il valore del Commodore non consisteva nei suoi circuiti, nella tastiera o nel monitor. Tutto questo era soltanto lo
strumento attraverso il quale un'idea prendeva forma. La vera magia nasceva quando un essere umano riusciva a
trasformare un pensiero in una sequenza di comandi comprensibili alla macchina.

Il professor Rossi Brunori lasciò che osservassimo lo schermo ancora per qualche secondo.

$\text{\textgreek{«}}$Vi sembra poco?$\text{\textgreek{»}}$

Qualcuno accennò un sorriso.

$\text{\textgreek{«}}$Capisco cosa state pensando. Abbiamo scritto una riga di programma e il computer ha semplicemente
visualizzato un saluto. Non sembra un risultato eccezionale.$\text{\textgreek{»}}$

Si fermò e appoggiò una mano sulla tastiera.

$\text{\textgreek{«}}$Ma provate a cambiare punto di vista. Fino a pochi minuti fa questa macchina non sapeva nemmeno
che esistesse la parola {\textquotedbl}CIAO{\textquotedbl}. Adesso l'ha scritta perché qualcuno glielo ha
insegnato.$\text{\textgreek{»}}$

Prese una breve pausa.

$\text{\textgreek{«}}$Ogni videogioco che avete visto, ogni programma che un giorno utilizzerete, ogni applicazione che
esisterà in futuro nascerà esattamente nello stesso modo. All'inizio c'è sempre una riga di codice. Poi un'altra. Poi
un'altra ancora. La differenza tra un programma semplice e uno straordinario non è il modo in cui comincia. È la
pazienza con cui qualcuno continua a costruirlo.$\text{\textgreek{»}}$

Quelle ultime parole caddero nell'aula con una naturalezza disarmante.

Il professore probabilmente le pronunciò pensando soltanto alla lezione di quel pomeriggio.

Io, invece, senza esserne ancora consapevole, avevo appena ricevuto una delle lezioni che avrebbero accompagnato tutta
la mia vita.

Il professor Rossi Brunori non cancellò immediatamente il programma.

Lasciò che quella semplice parola rimanesse sullo schermo ancora qualche istante, quasi volesse concederci il tempo
necessario per comprenderne il significato. Nessuno sembrava avere fretta di proseguire la lezione. In fondo, quel
piccolo esperimento aveva già raggiunto il suo scopo: ci aveva dimostrato che il computer non era un oggetto
misterioso, ma una macchina con la quale era possibile instaurare un dialogo, purché si imparasse a utilizzare il suo
linguaggio.

$\text{\textgreek{«}}$Adesso,$\text{\textgreek{»}}$ disse con tono tranquillo, $\text{\textgreek{«}}$vorrei farvi una
domanda.$\text{\textgreek{»}}$

Si voltò verso la classe, intrecciando le mani dietro la schiena.

$\text{\textgreek{«}}$Secondo voi, il programma che abbiamo scritto è intelligente?$\text{\textgreek{»}}$

Qualcuno sorrise.

La domanda sembrava quasi una provocazione.

$\text{\textgreek{«}}$No...$\text{\textgreek{»}}$ rispose una ragazza. $\text{\textgreek{«}}$Ha fatto soltanto quello
che gli abbiamo detto.$\text{\textgreek{»}}$

$\text{\textgreek{«}}$Esatto.$\text{\textgreek{»}}$

Il professore annuì lentamente.

$\text{\textgreek{«}}$E questa non è una sua debolezza. È la sua forza.$\text{\textgreek{»}}$

Notò il nostro sguardo interrogativo e continuò.

$\text{\textgreek{«}}$Molte persone pensano che un computer debba essere intelligente per essere utile. In realtà è vero
quasi il contrario. Un computer è utile proprio perché esegue sempre le stesse istruzioni nello stesso identico modo.
Se oggi ottenete un risultato e domani eseguite lo stesso programma, il comportamento sarà identico. È questa
affidabilità che rende possibile costruire applicazioni sempre più complesse.$\text{\textgreek{»}}$

Camminò lentamente tra i banchi, senza interrompere il filo del discorso.

$\text{\textgreek{«}}$Provate a immaginare un muratore che, ogni mattina, decidesse di cambiare la lunghezza del metro
con cui misura i mattoni. Oppure un falegname che stabilisse ogni giorno che un centimetro ha una misura diversa.
Nessuno riuscirebbe più a costruire nulla.$\text{\textgreek{»}}$

Qualche risata attraversò l'aula.

$\text{\textgreek{«}}$Il computer, invece, è ostinato. Se un'istruzione significa una certa cosa oggi, significherà la
stessa cosa anche domani, il mese prossimo e l'anno prossimo. È questo che permette ai programmatori di fidarsi delle
macchine.$\text{\textgreek{»}}$

Mentre parlava, mi rendevo conto che stava demolendo un'idea che avevo sempre dato per scontata. Avevo immaginato il
computer come qualcosa di straordinariamente complicato, quasi dotato di una propria volontà. Invece il professore
continuava a descriverlo come uno strumento estremamente disciplinato, capace di fare soltanto ciò che gli veniva
richiesto, ma di farlo con una precisione assoluta.

$\text{\textgreek{«}}$Per questo motivo,$\text{\textgreek{»}}$ proseguì tornando vicino alla cattedra,
$\text{\textgreek{«}}$non dovete mai prendervela con il computer quando un programma non
funziona.$\text{\textgreek{»}}$

Ci guardò uno per uno.

$\text{\textgreek{«}}$Almeno... quasi mai.$\text{\textgreek{»}}$

L'aula scoppiò a ridere.

$\text{\textgreek{«}}$Può capitare che una macchina si guasti, naturalmente. Ma nella maggior parte dei casi il computer
sta semplicemente eseguendo le istruzioni che qualcuno gli ha scritto. Se il risultato non è quello che vi aspettavate,
il primo sospettato non deve essere il Commodore.$\text{\textgreek{»}}$

Si batté una mano sul petto.

$\text{\textgreek{«}}$Il primo sospettato siete voi.$\text{\textgreek{»}}$

Le risate aumentarono ancora.

$\text{\textgreek{«}}$Non vi offendete. Vale anche per me. Dopo tanti anni continuo a scrivere programmi che ogni tanto
fanno qualcosa di diverso da ciò che avevo immaginato. Quando succede, invece di arrabbiarmi con il computer, provo a
chiedermi quale istruzione gli abbia dato e perché l'abbia interpretata proprio in quel modo.$\text{\textgreek{»}}$

Quelle parole suscitarono in me una riflessione che allora non ero ancora in grado di formulare con chiarezza. Fino a
quel momento avevo sempre pensato agli errori come a qualcosa da evitare. Il professore, invece, li trattava come
strumenti di lavoro. Non sembrava minimamente infastidito dall'idea di sbagliare. Al contrario, dava quasi
l'impressione che ogni errore rappresentasse un'occasione per capire meglio il comportamento della macchina.

Ritengo che quella sia stata una delle lezioni più preziose di tutto il corso.

Negli anni successivi avrei trascorso migliaia di ore davanti a un computer, scrivendo programmi, correggendo difetti,
cercando la causa di comportamenti inattesi. Ho imparato linguaggi molto diversi dal BASIC, ho lavorato con sistemi
infinitamente più potenti del Commodore 64 e ho affrontato problemi che, da bambino, non avrei nemmeno saputo
immaginare. Eppure, ogni volta che un programma non si comporta come previsto, mi torna ancora in mente quella frase
pronunciata con assoluta semplicità in un'aula di scuola media.

\emph{{\textquotedbl}Il primo sospettato siete voi.{\textquotedbl}}

Con il tempo ho scoperto che non era soltanto una buona regola per programmare.

Era un ottimo modo di imparare qualsiasi mestiere.

La settimana trascorse con una lentezza insolita.

A quell'età il tempo sembrava avere un ritmo tutto suo. Bastava aspettare una domenica, una gita scolastica o un
compleanno perché i giorni sembrassero non finire mai. Quella volta, però, ad attendermi non c'era un giocattolo nuovo
né una festa con gli amici. Ad aspettarmi c'era un'aula di scuola media, un professore e un Commodore 64.

Ripensavo continuamente a quella parola comparsa sullo schermo.

READY.

Era sorprendente come cinque lettere potessero continuare a occupare i miei pensieri per giorni interi. Cercavo di
ricordare ogni dettaglio della lezione: il colore dello schermo, il rumore dei tasti, il modo in cui il professore si
muoveva davanti alla tastiera e, soprattutto, la semplicità con cui era riuscito a convincerci che anche noi, un
giorno, saremmo stati in grado di scrivere programmi.

Non avevo ancora un computer a casa sul quale fare esperimenti. Tutto ciò che possedevo erano i ricordi della lezione
precedente e una quantità enorme di domande che continuavano ad affollarmi la mente.

Che cosa si poteva fare, oltre a scrivere una parola sullo schermo?

Come facevano i videogiochi a muovere i personaggi?

Chi aveva deciso che bastasse una parola come PRINT per convincere il computer a visualizzare un testo?

E soprattutto, se era possibile insegnare una cosa così semplice alla macchina, fino a dove ci si poteva spingere?

Erano domande ingenue, ma autentiche. Oggi so che molti di quei quesiti avrebbero richiesto anni di studio per ricevere
una risposta completa. All'epoca, invece, mi bastava sapere che qualcuno, in quell'aula, sembrava conoscere tutte le
risposte.

Quando arrivò finalmente il giorno della seconda lezione, entrai nella scuola con un'emozione completamente diversa
rispetto alla settimana precedente. La prima volta ero stato soprattutto curioso. Adesso, invece, avevo la sensazione
di tornare in un luogo familiare. Conoscevo già il percorso fino all'aula, sapevo dove il professore avrebbe sistemato
il Commodore e immaginavo persino il momento in cui avrebbe acceso il monitor.

Il professor Rossi Brunori era già arrivato.

Insieme a suo figlio stava terminando di collegare il monitor, il drive e il computer, controllando con calma che ogni
cavo fosse al proprio posto. Era un'operazione che eseguivano con una naturalezza quasi rituale, frutto di un'abitudine
consolidata. Noi, nel frattempo, prendevamo posto senza bisogno che nessuno ci invitasse a farlo. Era come se ciascuno
avesse ormai trovato spontaneamente il proprio ruolo all'interno di quel piccolo laboratorio improvvisato.

Non appena lo schermo si illuminò, il professore si voltò verso di noi.

$\text{\textgreek{«}}$Allora, vediamo un po'. Qualcuno si ricorda come si chiama il comando che abbiamo usato la
settimana scorsa?$\text{\textgreek{»}}$

Questa volta le mani si alzarono molto più rapidamente.

$\text{\textgreek{«}}$PRINT!$\text{\textgreek{»}}$ rispondemmo quasi all'unisono.

$\text{\textgreek{«}}$Ottimo.$\text{\textgreek{»}}$

Il professore annuì con soddisfazione.

$\text{\textgreek{«}}$Questo significa che non avete imparato una parola. Avete imparato un'idea. E le idee, quando sono
comprese davvero, rimangono molto più a lungo della memoria.$\text{\textgreek{»}}$

Senza aggiungere altro, digitò il piccolo programma della lezione precedente.

10 PRINT {\textquotedbl}CIAO{\textquotedbl}

Premette \textbf{RETURN}, poi scrisse il comando RUN.

La parola \textbf{CIAO} ricomparve sul monitor esattamente come la settimana precedente.

$\text{\textgreek{«}}$Bene,$\text{\textgreek{»}}$ disse sorridendo. $\text{\textgreek{«}}$La settimana scorsa abbiamo
insegnato al Commodore a salutarci una volta sola.$\text{\textgreek{»}}$

Si fermò un istante, lasciando che osservassimo il risultato.

$\text{\textgreek{«}}$Adesso immaginate di voler costruire un'insegna luminosa che continui a salutare tutte le persone
che passano davanti a un negozio. Oppure il tabellone di una stazione ferroviaria, che deve mostrare continuamente
un'informazione. Secondo voi, dovremmo tornare ogni volta alla tastiera per riscrivere il
programma?$\text{\textgreek{»}}$

Scuotemmo tutti la testa.

$\text{\textgreek{«}}$Naturalmente no.$\text{\textgreek{»}}$

Il professore si avvicinò al monitor e indicò la parola appena visualizzata.

$\text{\textgreek{«}}$Un programmatore, quando si accorge di dover ripetere la stessa operazione molte volte, cerca
sempre un modo per far lavorare il computer al posto suo. Anzi, vi dirò una cosa che forse vi
sorprenderà.$\text{\textgreek{»}}$

Fece una breve pausa, come era solito fare quando stava per consegnarci un'idea destinata a rimanere impressa.

$\text{\textgreek{«}}$Se vi trovate a svolgere continuamente la stessa operazione, probabilmente non siete voi a dover
lavorare di più. È il programma che deve diventare più intelligente.$\text{\textgreek{»}}$

Quella frase sembrò attraversare l'aula in un silenzio quasi assoluto.

Il professor Rossi Brunori sorrise.

$\text{\textgreek{«}}$Ed è proprio quello che proveremo a fare oggi.$\text{\textgreek{»}}$

Il professor Rossi Brunori rimase per qualche istante davanti al monitor senza toccare la tastiera. Era evidente che non
volesse limitarci a mostrare un nuovo comando. Come nelle lezioni precedenti, cercava prima di far nascere un'esigenza,
perché aveva capito che una soluzione viene ricordata molto più facilmente quando risponde a un problema che abbiamo
già percepito come nostro.

$\text{\textgreek{«}}$Vi faccio una proposta,$\text{\textgreek{»}}$ disse con il suo consueto sorriso.
$\text{\textgreek{«}}$Per qualche minuto dimenticate di essere davanti a un computer.$\text{\textgreek{»}}$

Si allontanò dalla cattedra e iniziò a camminare lentamente tra i banchi.

$\text{\textgreek{«}}$Immaginate di essere all'ingresso di una scuola e di dover salutare tutti gli studenti che
entrano. Arriva il primo.$\text{\textgreek{»}}$

Fece un cenno con la mano.

$\text{\textgreek{«}}$Ciao.$\text{\textgreek{»}}$

Ne fece un altro.

$\text{\textgreek{«}}$Arriva il secondo. Ciao.$\text{\textgreek{»}}$

Poi un terzo.

$\text{\textgreek{«}}$Ciao.$\text{\textgreek{»}}$

Si fermò.

$\text{\textgreek{«}}$Secondo voi, dopo il decimo studente, che cosa comincereste a pensare?$\text{\textgreek{»}}$

Qualcuno rise.

$\text{\textgreek{«}}$Che mi sono stancato.$\text{\textgreek{»}}$

$\text{\textgreek{«}}$Esatto.$\text{\textgreek{»}}$

Il professore annuì.

$\text{\textgreek{«}}$E dopo il centesimo?$\text{\textgreek{»}}$

$\text{\textgreek{«}}$Molto stancato.$\text{\textgreek{»}}$

Le risate riempirono nuovamente l'aula.

$\text{\textgreek{«}}$Vedete? Le persone si annoiano quando devono ripetere continuamente la stessa azione. I computer,
invece, no. Se una cosa deve essere eseguita cento volte, mille volte o un milione di volte, non si lamentano, non
perdono la concentrazione e non chiedono una pausa. Per loro ogni ripetizione è identica alla
precedente.$\text{\textgreek{»}}$

Tornò davanti al Commodore.

$\text{\textgreek{«}}$Ed è proprio per questo che esistono.$\text{\textgreek{»}}$

Quelle quattro parole, pronunciate con estrema semplicità, racchiudevano una verità che allora non ero in grado di
cogliere fino in fondo. Col passare degli anni avrei scoperto che una parte enorme del lavoro di un programmatore
consiste proprio nell'individuare attività ripetitive e trovare il modo di affidarle a una macchina. Non per sostituire
l'intelligenza dell'uomo, ma per liberarla da ciò che può essere automatizzato.

Il professore indicò nuovamente il programma.

$\text{\textgreek{«}}$Adesso il Commodore sa salutarci una sola volta.$\text{\textgreek{»}}$

Prese il gesso, lo rigirò tra le dita come faceva spesso quando rifletteva ad alta voce, ma ancora una volta lo ripose
senza utilizzarlo.

$\text{\textgreek{«}}$Secondo voi, come potremmo convincerlo a ricominciare da capo?$\text{\textgreek{»}}$

Nessuno parlò.

La domanda sembrava più difficile delle precedenti.

$\text{\textgreek{«}}$Proviamo a ragionare.$\text{\textgreek{»}}$

Indicò la prima riga.

$\text{\textgreek{«}}$Quando eseguiamo il programma, il computer legge la linea dieci, esegue il comando PRINT e scrive
la parola {\textquotedbl}CIAO{\textquotedbl}.$\text{\textgreek{»}}$

Con il dito seguì idealmente il percorso del programma.

$\text{\textgreek{«}}$E dopo?$\text{\textgreek{»}}$

Silenzio.

$\text{\textgreek{«}}$Che cosa trova dopo la linea dieci?$\text{\textgreek{»}}$

Uno dei ragazzi rispose quasi sottovoce.

$\text{\textgreek{«}}$Niente.$\text{\textgreek{»}}$

$\text{\textgreek{«}}$Esattamente.$\text{\textgreek{»}}$

Il professore sorrise.

$\text{\textgreek{«}}$Ed è proprio questo il motivo per cui il programma termina. Il Commodore non è pigro.
Semplicemente non sa più che cosa fare. Ha eseguito tutte le istruzioni che aveva a disposizione.$\text{\textgreek{»}}$

Rimase qualche secondo in silenzio.

$\text{\textgreek{«}}$Se voi steste leggendo un libro e, arrivati all'ultima pagina, non ce ne fossero altre, che cosa
fareste?$\text{\textgreek{»}}$

$\text{\textgreek{«}}$Lo chiuderei.$\text{\textgreek{»}}$

$\text{\textgreek{«}}$Certo.$\text{\textgreek{»}}$

Il professore annuì.

$\text{\textgreek{«}}$Il Commodore fa esattamente la stessa cosa. Quando non trova altre istruzioni, considera concluso
il proprio lavoro.$\text{\textgreek{»}}$

Si voltò verso lo schermo.

$\text{\textgreek{«}}$Ma noi possiamo anche dirgli una cosa diversa.$\text{\textgreek{»}}$

Le sue dita si posarono nuovamente sulla tastiera.

Per qualche istante nessuno nell'aula distolse lo sguardo dal monitor.

Il professore iniziò a scrivere lentamente una nuova riga.

20 GOTO 10

Premette \textbf{RETURN} e lasciò che osservassimo quelle poche parole.

$\text{\textgreek{«}}$Leggiamola insieme.$\text{\textgreek{»}}$

Attese che tutti fissassimo lo schermo.

$\text{\textgreek{«}}$Linea venti.$\text{\textgreek{»}}$

Indicò la seconda parola.

$\text{\textgreek{«}}$GOTO.$\text{\textgreek{»}}$

La pronunciò lentamente, separando quasi le due sillabe.

$\text{\textgreek{«}}$GO... TO...$\text{\textgreek{»}}$

Poi sorrise.

$\text{\textgreek{«}}$Non è una parola inventata dagli informatici. È semplicemente inglese e significa
{\textquotedbl}vai a{\textquotedbl}.$\text{\textgreek{»}}$

Indicò il numero che seguiva.

$\text{\textgreek{«}}$Questa istruzione dice al computer una cosa molto semplice: quando arrivi qui, invece di fermarti,
torna alla linea dieci.$\text{\textgreek{»}}$

Fece un piccolo gesto con la mano, tracciando nell'aria una specie di cerchio.

$\text{\textgreek{«}}$Dieci... venti... di nuovo dieci... poi ancora venti... e ancora dieci.$\text{\textgreek{»}}$

Continuava a muovere la mano seguendo quel percorso immaginario, finché uno dei ragazzi esclamò quasi d'istinto:

$\text{\textgreek{«}}$Ma allora non finisce più!$\text{\textgreek{»}}$

Il professore lo guardò con un'espressione soddisfatta.

Non rispose immediatamente.

Lasciò che quella frase rimanesse sospesa nell'aria, perché aveva capito che, questa volta, la scoperta non era nata da
una sua spiegazione.

Era nata dal ragionamento di uno studente.

Ed era proprio questo il risultato che cercava ogni volta che entrava in quell'aula.

Il professor Rossi Brunori si voltò lentamente verso il ragazzo che aveva pronunciato quelle parole. Sul suo volto
comparve un'espressione di sincera soddisfazione, la stessa che mostrava ogni volta in cui uno di noi arrivava a una
conclusione con le proprie forze.

$\text{\textgreek{«}}$Esatto,$\text{\textgreek{»}}$ disse con calma. $\text{\textgreek{«}}$Hai appena scoperto uno dei
concetti più importanti di tutta la programmazione.$\text{\textgreek{»}}$

Non aggiunse altro per qualche secondo. Non voleva che quella frase passasse inosservata.

$\text{\textgreek{«}}$Un programma non deve necessariamente avere una fine. Può anche continuare a lavorare finché
qualcuno non gli dice di fermarsi.$\text{\textgreek{»}}$

Lasciò che quell'idea trovasse posto nella nostra immaginazione.

$\text{\textgreek{«}}$Pensate all'orologio digitale che avete sul comodino di casa. Secondo voi, visualizza l'ora una
sola volta e poi smette?$\text{\textgreek{»}}$

$\text{\textgreek{«}}$No.$\text{\textgreek{»}}$

$\text{\textgreek{«}}$Oppure al semaforo che trovate ogni mattina andando a scuola. Cambia colore una sola volta nella
sua vita?$\text{\textgreek{»}}$

Qualcuno sorrise.

$\text{\textgreek{«}}$No.$\text{\textgreek{»}}$

$\text{\textgreek{«}}$E il tabellone delle partenze in una stazione? O l'insegna luminosa di una farmacia? O ancora il
registratore di cassa di un supermercato? Tutte queste macchine fanno la stessa cosa per ore, per giorni, a volte per
anni. Continuano a eseguire le proprie istruzioni senza stancarsi mai.$\text{\textgreek{»}}$

Si fermò davanti al Commodore e gli diede una leggera pacca sulla scocca.

$\text{\textgreek{«}}$Anche questo piccolo computer è stato costruito per fare esattamente la stessa
cosa.$\text{\textgreek{»}}$

Per la prima volta il Commodore 64 smetteva di sembrarmi un oggetto destinato esclusivamente ai videogiochi. Fino a quel
momento avevo sempre immaginato che un computer servisse soprattutto a divertirsi. Le parole del professore stavano
lentamente modificando quella convinzione. Ogni esempio che proponeva apparteneva alla vita quotidiana: un orologio, un
semaforo, un registratore di cassa. Senza rendercene conto, ci stava mostrando che l'informatica era già presente
intorno a noi, anche se non avevamo ancora imparato a riconoscerla.

$\text{\textgreek{«}}$Adesso, però, dobbiamo stare molto attenti.$\text{\textgreek{»}}$

Il tono della sua voce cambiò leggermente.

Non era diventato severo. Era il tono di chi sta per mettere in guardia degli apprendisti da un errore molto comune.

$\text{\textgreek{«}}$Vi ricordate che cosa vi ho detto la settimana scorsa? Il computer non è intelligente. Fa
esattamente ciò che gli chiediamo di fare.$\text{\textgreek{»}}$

Indicò nuovamente la riga appena scritta.

$\text{\textgreek{«}}$Se gli ordiniamo di tornare alla linea dieci, lui tornerà alla linea dieci. Se gli ordiniamo di
rifarlo ancora, lo rifarà. E continuerà così fino a quando qualcuno non interromperà il
programma.$\text{\textgreek{»}}$

Prese una breve pausa.

$\text{\textgreek{«}}$Non si fermerà perché pensa che sia abbastanza. Non si fermerà perché si annoia. Non si fermerà
perché capisce che il lavoro è finito.$\text{\textgreek{»}}$

Sorrise.

$\text{\textgreek{«}}$Non si fermerà nemmeno perché è ora di cena.$\text{\textgreek{»}}$

L'aula scoppiò a ridere.

$\text{\textgreek{«}}$Farà soltanto quello che gli avete ordinato.$\text{\textgreek{»}}$

Ancora una volta il professore era riuscito a trasformare un concetto tecnico in qualcosa di estremamente concreto. Oggi
definiremmo quel comportamento un ciclo infinito. Lui, invece, preferiva raccontarlo attraverso esempi che un gruppo di
ragazzi delle scuole elementari e medie poteva comprendere senza alcuna difficoltà.

Poi appoggiò entrambe le mani sulla tastiera.

$\text{\textgreek{«}}$Bene.$\text{\textgreek{»}}$

Pronunciò quella parola quasi sottovoce.

$\text{\textgreek{«}}$Adesso il nostro programma è davvero completo.$\text{\textgreek{»}}$

Sul monitor comparivano soltanto due righe.

10 PRINT {\textquotedbl}CIAO{\textquotedbl}

20 GOTO 10

Le osservammo in silenzio.

Erano incredibilmente brevi.

A vederle oggi sembrerebbero quasi banali. Allora, invece, rappresentavano qualcosa di sorprendente. Con appena due
istruzioni avevamo costruito un comportamento. Non avevamo più semplicemente chiesto al computer di scrivere una
parola. Gli avevamo insegnato che cosa fare dopo averla scritta.

Era una differenza enorme.

Il professor Rossi Brunori si voltò ancora una volta verso la classe.

$\text{\textgreek{«}}$Prima di eseguirlo, voglio chiedervi un'ultima cosa.$\text{\textgreek{»}}$

Ormai conoscevamo bene quel rituale. Ogni dimostrazione era preceduta da una domanda, perché il professore voleva che la
risposta nascesse nella nostra mente prima ancora di apparire sul monitor.

$\text{\textgreek{«}}$Questa volta non vi chiedo che cosa farà il programma. Lo sappiamo già.$\text{\textgreek{»}}$

Fece un piccolo sorriso.

$\text{\textgreek{«}}$Vi chiedo una cosa diversa. Se davvero continuerà a scrivere {\textquotedbl}CIAO{\textquotedbl}
senza fermarsi mai, secondo voi, come faremo a interromperlo?$\text{\textgreek{»}}$

Nell'aula tornò il silenzio.

Nessuno aveva mai pensato che potesse esistere anche il problema opposto. Fino a pochi minuti prima ci stavamo
preoccupando di convincere il computer a fare qualcosa. Adesso, improvvisamente, il problema era capire come impedirgli
di continuare.

Fu in quel momento che compresi, senza ancora esserne pienamente consapevole, una caratteristica che avrei ritrovato in
tutta l'informatica degli anni successivi.

Ogni soluzione genera inevitabilmente una nuova domanda.

Ed è proprio quella domanda ad aprire la strada alla scoperta successiva.

Il professor Rossi Brunori osservò ancora una volta il piccolo programma sullo schermo. Due sole righe. A guardarle con
gli occhi di un adulto potevano sembrare quasi insignificanti, eppure in quell'aula rappresentavano qualcosa di molto
più grande della loro apparente semplicità.

$\text{\textgreek{«}}$Vediamo se abbiamo ragionato bene.$\text{\textgreek{»}}$

Premette il tasto \textbf{RETURN} dopo il comando RUN.

Per una frazione di secondo il monitor rimase immobile. Poi la parola \textbf{CIAO} cominciò a scendere sullo schermo
con una velocità che nessuno di noi si aspettava.

CIAO

CIAO

CIAO

CIAO

CIAO

Le righe si susseguivano senza interruzione, sempre più numerose, fino a riempire completamente il monitor. Quando lo
schermo arrivava in fondo, il testo ricominciava a scorrere dall'alto in un movimento continuo che sembrava non avere
fine.

Nell'aula si levò un brusio di stupore.

Qualcuno rise.

Qualcuno si alzò leggermente dalla sedia per osservare meglio.

Altri cercavano inutilmente di leggere tutte quelle scritte che scorrevano troppo in fretta.

Il Commodore sembrava improvvisamente diventato vivo.

Non stava più aspettando istruzioni.

Le stava eseguendo senza esitazione, con una velocità che ai nostri occhi era quasi incredibile.

Il professore lasciò che quella dimostrazione continuasse ancora qualche secondo. Poi allungò una mano e premette un
tasto sulla tastiera.

Lo schermo si fermò immediatamente.

Il silenzio che seguì fu quasi sorprendente quanto il programma stesso.

$\text{\textgreek{«}}$Questo tasto si chiama RUN/STOP.$\text{\textgreek{»}}$

Lo indicò senza alcuna enfasi.

$\text{\textgreek{«}}$Ricordatevelo. Non perché dovrete usarlo spesso, ma perché vi insegna una cosa
importante.$\text{\textgreek{»}}$

Attese che tutti lo guardassimo.

$\text{\textgreek{«}}$Quando scrivete un programma, dovete sempre sapere come riprenderne il controllo. Un computer
esegue gli ordini con una precisione straordinaria, ma la responsabilità di quegli ordini rimane sempre di chi li ha
scritti.$\text{\textgreek{»}}$

Quelle parole mi tornarono in mente moltissime volte negli anni successivi. Ogni volta che un programma entrava in un
ciclo inatteso, ogni volta che un processo sembrava non terminare mai, ogni volta che un'applicazione smetteva di
rispondere, riaffiorava il ricordo di quella semplice dimostrazione fatta davanti a un Commodore 64.

Non era una lezione sul tasto RUN/STOP.

Era una lezione sulla responsabilità del programmatore.

Il professor Rossi Brunori spense il monitor.

Per noi quello era il segnale che la lezione era finita.

Cominciammo lentamente a raccogliere quaderni e penne, continuando a commentare tra noi quello che avevamo appena visto.
Le conversazioni erano tutte diverse, ma avevano un elemento in comune: nessuno parlava di compiti, di interrogazioni o
di voti. Parlavammo soltanto del computer.

Mentre uscivo dall'aula mi voltai ancora una volta.

Il Commodore 64 era ormai spento. Sembrava un oggetto qualunque, quasi anonimo, appoggiato sulla cattedra. Era difficile
credere che pochi minuti prima fosse riuscito a catturare completamente la nostra attenzione con appena due righe di
programma.

Solo molti anni dopo avrei capito che non era stato il computer ad affascinarmi.

Era l'idea.

L'idea che una macchina potesse fare esattamente ciò che avevo immaginato, purché fossi capace di spiegarglielo nel modo
giusto.

Fu quella convinzione, più ancora del Commodore stesso, ad accompagnarmi sulla strada di casa.

E, senza che potessi saperlo, era destinata a cambiare il corso della mia vita.


\bigskip
\end{document}
